\section{Background (Draft v2)}
\label{sec:background}

% \todo{Introduce stuff necessary for following sections. e.g. wan 2.1 vs 2.2, video generation, latent space, rope etc @ADI Working on it}

\subsection{Denoising-based Video Diffusion Models}

Recent video diffusion models increasingly uses flow matching~\cite{flow_matching_lipman2022flow}, a type of stochastic-interpolant~\cite{stochastic_interpolants} as the objective  and diffusion transformers ~\cite{peebles2023scalable,flow_matching_ma2024sit} as model backbones for denoising.

Following Rectified Flows (RFs) ~\cite{flow_matching_liu2022flow}, video diffusion models learn the velocity along a linear path from noise $x_0\!\sim\!\mathcal{N}(0,I)$ to data $x_1$. For a sampled time $t\in[0,1]$ the intermediate latent is $x_t = t x_1 + (1-t)x_0$ and the ground-truth velocity is $v_t = \mathrm{d}x_t/\mathrm{d}t = x_1 - x_0$. The model $u(\cdot;\theta)$ predicts this velocity conditioned on $x_t$, $t$ and context $c_{\text{context}}$ (e.g., text, image or camera parameters) by minimizing the MSE objective-
\[
\mathcal{L} = \mathbb{E}_{x_0,x_1,c_{\text{context}},t}\big[\|u(x_t,c_{\text{context}},t;\theta)-v_t\|^2\big].
\]

This formulation compactly casts generation as learning a deterministic velocity field along straight-line interpolants between noise and data, enabling conditional video synthesis with flow-matching objectives.

\subsection{DiT based conditional video diffusion models}

The advent of large-scale open-source Transformer-based (DiT) video diffusion models, like Wan2.1 ~\cite{wan21}, CogVideoX ~\cite{cogvideox} and HunyuanVideo ~\cite{kong2024hunyuanvideo}, has catalyzed the development of techniques to condition these models on diverse auxiliary inputs. The methodologies for integrating these conditional signals are multi-faceted. A common approach involves first encoding the control video or image into the latent space via the VAE and subsequently introducing these latent tokens as direct inputs to the diffusion model.

More direct integration strategies are also prevalent, often falling into three primary categories: (1) concatenation of the control signal with the noisy latent along the channel dimension; (2) injection of conditional information through dedicated cross-attention layers; or (3) introduction of control signals as supplementary tokens processed by the full self-attention mechanism.

For instance, \recammaster{}~\cite{bai2025recammaster} demonstrates that frame-level concatenation is a particularly robust method for integrating video controls within DiT-based architectures. Similarly, other papers like \trajectorycrafter{}~\cite{yu2025trajectorycrafter} and \exfd{}~\cite{ex4d} utilize concatenation along the channel dimension of the noisy latent, a technique that spatially aligns the control tokens with the nascent generation.

In contrast, models like VACE~\cite{jiang2025vace}, along with other subject-consistent video generation methods, often employ image conditioning in the form of tokens, with supplementary embeddings looped into the model through cross-attention layers.

The choice of conditioning mechanism has significant implications, especially regarding temporal modeling and positional embeddings. Rotary Positional Embeddings ~\cite{su2024roformer}, for example, play a critical role. Methods such as Ovi ~\cite{low2025ovi} leverage RoPE aligned with audio inputs to achieve robust speech-video-audio alignment. Conversely, approaches like \recammaster{}~\cite{bai2025recammaster}, which append additional tokens at the end of the sequence, can be limited to producing only a fixed, pre-defined video length, highlighting the crucial interplay between the conditioning strategy and the model's temporal architecture.


\todo{don't need repeat citations for near acronym reuse}